\section{The arithmetic certificate}\label{sec:arith-certificate}

This appendix documents the computation behind \cref{prop:A-value}, the exact value of the
arithmetic constant $\cA$ used in \cref{thm:arithmetic}, and the checks it has passed. It
describes how the type data are computed, how the profile is integrated exactly, and how the local
bound of \cref{prop:nullspace} was tested against exact valuations at finite $h$.

\subsection{Type data}\label{subsec:type-data}

For each $\kappa\in\{\frac15,\dots,\frac45\}$ and each type
$(\kappa,\mathcal L^\Row,\mathcal L^\Col,\varepsilon)$ of the alternating form of
\cref{def:profile} (common digits $\{0,\dots,\gamma-1\}$, border digits $\{\gamma,\dots,T-1\}$ alternating
between rows and columns, $T\le23$), the program computes the rational numbers $V_l,W_l$ of
\eqref{eq:VW-def}, the block $A_\tau$, its rank over $\Q$, bases of its left and right null spaces
by exact row reduction, the images of these null spaces under the moment map, and the lines they
span. Over the whole range $1/20\le u\le45/8$ about a thousand types occur ($1042$
according to the AI-internal audit), and no null space has a two-dimensional moment image, so the
convention of \cref{def:block-types}(ii) for two-dimensional images is never used.

\begin{example}\label{ex:profile-x4}
Let $\kappa=\frac15$ and $u=4$. For $t\in[0,1)$ a fibre has one common digit; for
$t\in[1,\frac98)$ one border digit; for larger $t$ none. Here $\varepsilon(t)=(\mathbf1[t>\frac45],0)$.
A common singleton has $A=[W_0]$ with $W_0=-\varepsilon_1/\kappa^2$: rank $0$ for $t<\frac45$, with
left and right null vector $(1)$ of moment $(1,V_0)=(1,0)$, and rank $1$ for $\frac45<t<1$. A border
singleton is a $1\times0$ or $0\times1$ block of rank $0$ whose null vector has moment
$(1,1/\kappa)=(1,5)$; half of these fibres are rows and half columns. Hence $\rho=0.2$,
$\alpha=\beta=0.8+0.0625$, the line $\Q(1,0)$ has $\alpha_\Lambda=\beta_\Lambda=0.8$ and the line
$\Q(1,5)$ has $\alpha_\Lambda=\beta_\Lambda=0.0625$. Thus $q=\min(0.8625,\,0.8625,\,0.125,\,1.6)=0.125$
and $e_{1/5}(4)=0.4+0.125=0.525$. Without the line bound one would get $0.4+0.8625$.
\end{example}

\Cref{tab:profile} shows the profile at a few points. It falls from about $2.1$, close to
the trivial value $2+c$, to $0$, but not monotonically (for $\kappa=\frac25$ it is $2.000$ near
$u=0.5$ and $2.050$ near $u=0.75$); for $\kappa\ge\frac25$ it vanishes for $u\ge3$, while for
$\kappa=\frac15$ it stays positive up to $u=45/8$, which explains the larger integral in that
class.

\begin{table}[ht]
\centering
\small
\setlength{\tabcolsep}{4.6pt}
\begin{tabular}{@{}lcccccccccccc@{}}
\toprule
$u$ & $0.05$ & $0.1$ & $0.25$ & $0.5$ & $0.75$ & $1$ & $1.5$ & $2$ & $3$ & $4$ & $5$ & $5.5$\\
\midrule
$\kappa=\frac15$ & 2.110 & 2.075 & 2.050 & 2.025 & 2.000 & 1.850 & 1.525 & 1.325 & 0.925 & 0.525 & 0.125 & 0.025\\
$\kappa=\frac25$ & 2.105 & 2.075 & 2.025 & 2.000 & 2.050 & 1.450 & 0.925 & 0.525 & 0 & 0 & 0 & 0\\
$\kappa=\frac35$ & 2.100 & 2.085 & 2.000 & 2.000 & 2.000 & 1.250 & 0.625 & 0.125 & 0 & 0 & 0 & 0\\
$\kappa=\frac45$ & 2.100 & 2.0875 & 2.000 & 2.000 & 2.000 & 1.250 & 0.625 & 0.125 & 0 & 0 & 0 & 0\\
\bottomrule
\end{tabular}
\caption{The profile $e_\kappa$ at the points $u+10^{-6}$ for the displayed $u$, rounded to three
decimals (the entry $2.0875$ is given with four). The trivial
entrywise bound corresponds to the constant $2+c=2.125$.}
\label{tab:profile}
\end{table}

\subsection{Exact integration}\label{subsec:exact-integration}

For each $\kappa$ the interval $[1/20,45/8]$ is cut at the points~\eqref{eq:breakpoints} lying
inside it, giving $122$, $123$, $122$ and $122$ cells for $\kappa=\frac15,\dots,\frac45$. On a cell
$(l,r)$ the original program evaluates all candidate functions
$2\rho+\alpha$, $2\rho+\beta$ and $2\rho+\alpha+\beta-\alpha_\Lambda-\beta_\Lambda$ at the interior
points $(2l+r)/3$ and $(l+2r)/3$, reconstructs each as an affine function, adds the pairwise
intersection points inside the cell as further cuts, and integrates the minimum exactly by
evaluating it at the midpoint of each sub-cell. Adding $\frac{2+c}{20}=\frac{17}{160}$ for the
range below $u=1/20$ and averaging over $\kappa$ gives
\[
  \cA=\frac{304227995343501721142691896310988016794079029050925423}
           {86622067750624596685118168422516024981216422248448000}.
\]
The inequality $\cA<3512131/10^6$ is checked in rational arithmetic. No quadrature or floating
point enters.

\subsection{Checks of the arithmetic computation}\label{subsec:arith-checks}

\begin{enumerate}[(a)]
  \item \emph{Separated geometry.} The original program integrates the analogous profile for the
  separated geometry (border rows on $[1,1+c/2)$, columns on $[1+c/2,1+c)$, no parity averaging)
  and obtains exactly the same rational number. This is consistent with \cref{cor:separation},
  which makes the odd-prime denominators of the two geometries equal.
  \item \emph{Recomputation from the description.} In the AI-internal audit the profile was
  reimplemented from the mathematical description, with the piece lengths carried symbolically
  as affine functions of $u$, the breakpoints generated independently as the intersections of
  the lines $0,u,\kappa u,u/2,1-ju,1+c-ju$, and the constancy of the combinatorial type and of the
  affine candidates asserted at two points of every interval. It reproduces the four integrals
  and $\cA$ exactly. Using the best labelling of two-dimensional moment images instead of the
  coordinate lines gives the same value, since such images do not occur.
  \item \emph{Affinity.} We evaluated, for each of the $489$ cells of the original program, the
  candidate functions at the three further points $l+\frac1{11}(r-l)$, $\frac{l+r}2$,
  $l+\frac9{10}(r-l)$ and compared them with the affine reconstruction: all agree exactly.
  Together with the two points of the reconstruction this makes five points per cell. Such a test
  supports \cref{lem:profile-structure} but cannot replace it: finitely many points do not
  exclude a breakpoint inside a cell.
  \item \emph{Polar congruences.} The congruences of \cref{lem:two-layer} were checked on
  $312{,}120$ matrix entries (\cref{subsec:polar}).
\end{enumerate}

\subsection{Finite tests of the local bound}\label{subsec:finite-checks}

\Cref{prop:nullspace} can be tested at finite $h$ by computing the exact polynomial
$P^{\alt}_I(X)$, or its value at a rational point, and comparing the exponent $e_p$ with the
prediction $2R_0+q$ computed from the actual fibres. Since $p_0$ has not been computed, we cannot
say which of the primes tested are ordinary; those below $120$, where the exponents are largest,
are not. The tests therefore check the formula of the proposition, and the mechanism of its proof,
partly at primes that the proof excludes. Recall from
\cref{rem:bad-primes} that a prime is \emph{bad} if some fibre block has smaller rank modulo $p$
than over $\Q$.
\begin{enumerate}[(a)]
  \item \emph{Rational data, value at a point.} We compared $e_p(P^{\alt}_\varnothing(1/7))$ with
  $2R_0+q$, computed over $\Q$ exactly as in \cref{prop:nullspace}, at every prime
  $p\notin\{5,7\}$ with $\sqrt{5N}<p\le5N$, for $h\in\{16,32,48,64,80,96,128\}$ and all $k$
  ($1836$ comparisons): no violation, at bad primes or elsewhere; equality at $1810$ primes, and
  at $117$ to $119$ of the $119$ primes when $h=128$. (The sizes $h\in\{32,48,64,96\}$ are run by
  a program of the package; the others were added for this revision.)
  \item \emph{Data modulo $p$, Gauss valuations.} The AI-internal audit compared the Gauss
  valuations of the exact polynomials $P^{\alt}_I(X)$ (\cref{rem:gauss}) with the number
  $2R'_0+\rank_{\Fp}K$ in which every rank is computed modulo $p$: $R'_0$ is the sum of the ranks
  over $\Fp$ of the fibre blocks, and $K$ is the pairing between their null spaces over $\Fp$.
  At a prime that is not bad, $R'_0=R_0$ and this is the intermediate bound in the proof of
  \cref{prop:nullspace}. The comparison was made at every odd prime $p\ne5$ with $p^2>5N$, for
  $h\in\{16,32,48,64,80,96,128\}$ and all $k$, and with the deletions $\{0\}$ and $\{h-1\}$ for
  $h\in\{32,48\}$: no violation at any prime that is not bad.
  \item \emph{Bad primes.} Four bad primes occur in the tests of (b) without deletion
  (\cref{tab:bad-primes}); the prime $17$ at $h=32$, $k=1$ is also bad for the two deletions. At
  two of the four, $e_p$ exceeds the number computed modulo $p$: by $3$ at $h=16$, $k=3$, $p=11$,
  and by $4$ at $h=128$, $k=3$, $p=31$, the largest size tested. As explained in
  \cref{rem:bad-primes}, this contradicts no statement of this paper: the number computed modulo
  $p$ is not the bound of \cref{prop:nullspace}, which is stated with ranks over $\Q$ and for
  $p>p_0$ only. With the rational data the inequality $e_p\le2R_0+q$ holds at all four bad primes,
  for the Gauss valuation and at the points $X\in\{0,1,\frac17,\frac23,\frac3{11}\}$. The rank
  drops are explicit: for $p=31$ the block of the type with four common digits, $\kappa=\frac35$ and
  $\varepsilon=(1,1)$ is a $4\times4$ matrix of rank $4$ whose determinant has numerator divisible
  by $31$, and this type occurs in the eleven fibres $20\le\iota\le30$; for $p=11$ and $p=13$ it
  is the $1\times1$ block $W_0=4-\kappa^{-2}$, equal to $\frac{11}9$ and $\frac{39}{16}$; for
  $p=17$ a $2\times2$ block of determinant $\frac{17}3$.
  \item \emph{Comparison with the profile.} In the tests of (a) the prediction $2R_0+q$ exceeds
  $h\,e_\kappa(p/h)$ by at most $0.8$, and the exponent $e_p$ itself exceeds $h\,e_\kappa(p/h)$ at
  $56$ of the $1836$ primes (for instance $98>97.2$ at $h=48$, $k=1$, $p=41$). So the profile does
  not dominate prime by prime; this is consistent with \cref{lem:profile-bound}, whose constant
  $C_0$ cannot be taken to be $0$. What holds in all $28$ cases is the inequality between the
  sums: $\sum_p(2R_0+q)\log p<\sum_ph\,e_\kappa(p/h)\log p$ over the primes tested, the difference
  being between $0.0099$ and $0.1754$ per $h^2$ and decreasing with $h$ for each $k$.
\end{enumerate}

\begin{table}[ht]
\centering
\small
\begin{tabular}{@{}ccccccc@{}}
\toprule
$h$ & $k$ & $p$ & $e_p$ & computed modulo $p$ & $2R_0+q$ over $\Q$ & $h\,e_\kappa(p/h)$\\
\midrule
$16$ & $2$ & $13$ & $28$ & $28$ & $28$ & $29.0$\\
$16$ & $3$ & $11$ & $29$ & $26$ & $33$ & $32.9$\\
$32$ & $1$ & $17$ & $64$ & $64$ & $66$ & $66.0$\\
$128$ & $3$ & $31$ & $253$ & $249$ & $260$ & $260.5$\\
\bottomrule
\end{tabular}
\caption{The four bad primes met in the finite tests (no deletion). Here $e_p$ is the exponent of
$p$ in the denominator of the content of $P^{\alt}_\varnothing$, and the fifth column is the
number $2R'_0+\rank_{\Fp}K$ of test~(b).}
\label{tab:bad-primes}
\end{table}

These tests concern the inequality of \cref{prop:nullspace}, not the asymptotic constant: at
these sizes the contribution of the small primes is still of the same order as that of the
others (\cref{subsec:margin}).
